\documentclass[10pt,a4paper]{article}

\usepackage{cmap}
\usepackage[utf8]{inputenc}
\usepackage[T1]{fontenc}
\usepackage{lmodern}
\usepackage[french]{babel}
\usepackage{array}
\usepackage{geometry}
\usepackage{longtable}
\usepackage{xcolor}
\usepackage{colortbl}
\usepackage{booktabs}
\usepackage{ragged2e}

\geometry{a4paper, landscape, margin=1.05cm}
\renewcommand{\arraystretch}{1.2}
\setlength{\tabcolsep}{4pt}

\definecolor{headergreen}{RGB}{169, 204, 120}
\definecolor{titleblue}{RGB}{0, 128, 170}
\definecolor{sectionpink}{RGB}{174, 37, 105}

\newcolumntype{P}[1]{>{\RaggedRight\arraybackslash}p{#1}}

\begin{document}

\begin{center}
    {\Large\bfseries\color{titleblue} Catégories de tests logiciels}\\[0.35cm]
    {\large\bfseries\color{sectionpink} Conception et gestion des cas de test}\\[0.2cm]
    {\large\bfseries Format de cas de test (modèle 1)}\\[0.35cm]
\end{center}

\noindent\textbf{Projet :} API de transaction bancaire\\
\textbf{Étudiant :} Demano Bramael\\
\textbf{Matricule :} 25G2033\\
\textbf{Fonctions testées :} création de compte, liste des comptes, consultation, dépôt et retrait.

\vspace{0.35cm}

{\small
\begin{longtable}{|P{2.1cm}|P{4.8cm}|P{7.1cm}|P{6.2cm}|P{2.4cm}|}
\hline
\rowcolor{headergreen}
\textbf{Test Case ID} &
\textbf{Test Case Description} &
\textbf{Test Input (Test Case Value)} &
\textbf{Expected Results} &
\textbf{Test Case Status} \\
\hline
\endfirsthead

\hline
\rowcolor{headergreen}
\textbf{Test Case ID} &
\textbf{Test Case Description} &
\textbf{Test Input (Test Case Value)} &
\textbf{Expected Results} &
\textbf{Test Case Status} \\
\hline
\endhead

TC-001 &
Créer un compte bancaire valide. &
\textbf{Méthode :} POST\newline
\textbf{Endpoint :} /api/comptes/creer/\newline
\textbf{Données :} numero\_compte = ACC000001, nom\_titulaire = DEMANO BRAMAEL, type\_compte = COURANT, solde = 10000.00 &
Code HTTP 201. Le compte est créé avec succès et une transaction de dépôt initial est enregistrée. &
Réussi \\
\hline

TC-002 &
Lister les comptes bancaires actifs. &
\textbf{Méthode :} GET\newline
\textbf{Endpoint :} /api/comptes/\newline
\textbf{Données :} aucune donnée dans le corps de la requête. &
Code HTTP 200. Le système retourne un tableau contenant les comptes actifs. &
Réussi \\
\hline

TC-003 &
Consulter les détails d'un compte bancaire. &
\textbf{Méthode :} GET\newline
\textbf{Endpoint :} /api/comptes/\{id\}/\newline
\textbf{Données :} id = 1 &
Code HTTP 200. Le système retourne les informations du compte et ses transactions. &
Réussi \\
\hline

TC-004 &
Effectuer un dépôt sur un compte actif. &
\textbf{Méthode :} POST\newline
\textbf{Endpoint :} /api/comptes/\{id\}/depot/\newline
\textbf{Données :} montant = 5000.00, description = Dépôt agence &
Code HTTP 200. Le solde augmente et une transaction de type DEPOT est enregistrée. &
Réussi \\
\hline

TC-005 &
Effectuer un retrait avec solde suffisant. &
\textbf{Méthode :} POST\newline
\textbf{Endpoint :} /api/comptes/\{id\}/retrait/\newline
\textbf{Données :} montant = 2000.00, description = Retrait guichet &
Code HTTP 200. Le solde diminue et une transaction de type RETRAIT est enregistrée. &
Réussi \\
\hline
\end{longtable}
}

\vspace{0.2cm}

\noindent\textbf{Conclusion :} les cinq fonctions principales du système bancaire ont été testées manuellement et les résultats attendus correspondent au comportement de l'API.

\end{document}
